\documentclass[10pt,a4paper]{amsart}
\usepackage{amsmath,amssymb,amsthm,graphicx}
\usepackage{verbatim}
\newenvironment{reference}{\comment}{\endcomment}
%\newenvironment{reference}{}{}
\newenvironment{slogan}{\comment}{\endcomment}
\newenvironment{history}{\comment}{\endcomment}
% For commutative diagrams we use Xy-pic
\usepackage[all]{xy}

% We use 2cell for 2-commutative diagrams.
\xyoption{2cell}
% Theorem environments.
%
\theoremstyle{plain}
\newtheorem{theorem}[subsection]{Theorem}
\newtheorem{proposition}[subsection]{Proposition}
\newtheorem{lemma}[subsection]{Lemma}

\theoremstyle{definition}
\newtheorem{definition}[subsection]{Definition}
\newtheorem{example}[subsection]{Example}
\newtheorem{exercise}[subsection]{Exercise}
\newtheorem{situation}[subsection]{Situation}

\theoremstyle{remark}
\newtheorem{remark}[subsection]{Remark}
\newtheorem{remarks}[subsection]{Remarks}

\numberwithin{equation}{subsection}

% Macros
%
\def\lim{\mathop{\mathrm{lim}}\nolimits}
\def\colim{\mathop{\mathrm{colim}}\nolimits}
\def\Spec{\mathop{\mathrm{Spec}}}
\def\Hom{\mathop{\mathrm{Hom}}\nolimits}
\def\Ext{\mathop{\mathrm{Ext}}\nolimits}
\def\SheafHom{\mathop{\mathcal{H}\!\mathit{om}}\nolimits}
\def\SheafExt{\mathop{\mathcal{E}\!\mathit{xt}}\nolimits}
\def\Sch{\mathit{Sch}}
\def\Mor{\mathop{\mathrm{Mor}}\nolimits}
\def\Ob{\mathop{\mathrm{Ob}}\nolimits}
\def\Sh{\mathop{\mathit{Sh}}\nolimits}
\def\NL{\mathop{N\!L}\nolimits}
\def\CH{\mathop{\mathrm{CH}}\nolimits}
\def\proetale{{pro\text{-}\acute{e}tale}}
\def\etale{{\acute{e}tale}}
\def\QCoh{\mathit{QCoh}}
\def\Ker{\mathop{\mathrm{Ker}}}
\def\Im{\mathop{\mathrm{Im}}}
\def\Coker{\mathop{\mathrm{Coker}}}
\def\Coim{\mathop{\mathrm{Coim}}}

% Boxtimes
%
\DeclareMathSymbol{\boxtimes}{\mathbin}{AMSa}{"02}

%
% Macros for moduli stacks/spaces
%
\def\QCohstack{\mathcal{QC}\!\mathit{oh}}
\def\Cohstack{\mathcal{C}\!\mathit{oh}}
\def\Spacesstack{\mathcal{S}\!\mathit{paces}}
\def\Quotfunctor{\mathrm{Quot}}
\def\Hilbfunctor{\mathrm{Hilb}}
\def\Curvesstack{\mathcal{C}\!\mathit{urves}}
\def\Polarizedstack{\mathcal{P}\!\mathit{olarized}}
\def\Complexesstack{\mathcal{C}\!\mathit{omplexes}}
% \Pic is the operator that assigns to X its picard group, usage \Pic(X)
% \Picardstack_{X/B} denotes the Picard stack of X over B
% \Picardfunctor_{X/B} denotes the Picard functor of X over B
\def\Pic{\mathop{\mathrm{Pic}}\nolimits}
\def\Picardstack{\mathcal{P}\!\mathit{ic}}
\def\Picardfunctor{\mathrm{Pic}}
\def\Deformationcategory{\mathcal{D}\!\mathit{ef}}
\usepackage{cleveref}
\crefname{originallemma}{original lemma}{original lemmas}
\Crefname{originallemma}{Original lemma}{Original lemmas}
\makeatletter
\newlabel{algebra-lemma-one-equation}{{00KW}{source}{}{sourcepointer-1}{}}
\newlabel{algebra-lemma-one-equation@cref}{{[originallemma][0][]00KW}{[1][0][]source}{}{}{}}
\makeatother
\crefname{originaldefinition}{original definition}{original definitions}
\Crefname{originaldefinition}{Original definition}{Original definitions}
\makeatletter
\newlabel{divisors-definition-Weil-divisor}{{0BE2}{source}{}{sourcepointer-2}{}}
\newlabel{divisors-definition-Weil-divisor@cref}{{[originaldefinition][0][]0BE2}{[1][0][]source}{}{}{}}
\makeatother
\crefname{originallemma}{original lemma}{original lemmas}
\Crefname{originallemma}{Original lemma}{Original lemmas}
\makeatletter
\newlabel{divisors-lemma-effective-Cartier-in-points}{{0AG8}{source}{}{sourcepointer-3}{}}
\newlabel{divisors-lemma-effective-Cartier-in-points@cref}{{[originallemma][0][]0AG8}{[1][0][]source}{}{}{}}
\makeatother
\crefname{originallemma}{original lemma}{original lemmas}
\Crefname{originallemma}{Original lemma}{Original lemmas}
\makeatletter
\newlabel{properties-lemma-codimension-local-ring}{{02IZ}{source}{}{sourcepointer-4}{}}
\newlabel{properties-lemma-codimension-local-ring@cref}{{[originallemma][0][]02IZ}{[1][0][]source}{}{}{}}
\makeatother
\crefname{originalreference}{original reference}{original references}
\Crefname{originalreference}{Original reference}{Original references}
\makeatletter
\newlabel{situation-setup}{{02QL}{source}{}{sourcepointer-5}{}}
\newlabel{situation-setup@cref}{{[originalreference][0][]02QL}{[1][0][]source}{}{}{}}
\makeatother
\crefname{originaldefinition}{original definition}{original definitions}
\Crefname{originaldefinition}{Original definition}{Original definitions}
\makeatletter
\newlabel{topology-definition-dimension-function}{{02I9}{source}{}{sourcepointer-6}{}}
\newlabel{topology-definition-dimension-function@cref}{{[originaldefinition][0][]02I9}{[1][0][]source}{}{}{}}
\makeatother
\title{Prime divisors and the components of an effective Cartier divisor in terms of the relative dimension function}
\author{The Stacks Project Authors}
\date{stacks-project:chow.tex:local-snapshot-20261012; GNU FDL 1.2 or later}
\begin{document}
\maketitle
\noindent\textit{Editorial note.} All delivered mathematical text is the author's own, copied from the pinned Stacks Project chapter file chow.tex; no statement or proof step was rewritten, reordered or removed.\par\medskip
\noindent\textit{Editorial note.} Preamble provenance: the pinned chapter file opens with \textbackslash{}input\{preamble\} and carries no preamble text of its own, so the project's single Stacks-wide preamble staged at sources/stacks-chow/preamble.tex is excerpted verbatim as three spans, lines 8-12 (the verbatim package together with the project's own reference, slogan and history display-environment definitions, including the author's commented-out older form of the reference environment), lines 14-18 (Xy-pic setup) and lines 154-217 (the native theorem declarations and mathematical macros). The lines 8-12 span is required because a counted statement or its surrounding text may use those display environments, which the project typesets as comments. Nothing was added, removed or redefined; the wrapper class and the amsmath/amssymb/amsthm/graphicx base packages are supplied by the assembler.\par\medskip
\noindent\textit{Editorial note.} Same-chapter pointers are located at their \textbackslash{}label lines inside the pinned chapter file; cross-chapter pointers are located in the project's own tag list staged in the snapshot (sources/stacks-chow/tags.txt), at the verbatim lines carrying their canonical tags.\par\medskip
\noindent\textit{Allowed standard prerequisites (editorial).}\begin{itemize}
\item The definition of a prime divisor, of a dimension function, and the local shape of an effective Cartier divisor. Divisors, labels divisors-definition-Weil-divisor (proof line 2508), divisors-lemma-effective-Cartier-in-points (proof line 2515); Topology, label topology-definition-dimension-function (proof line 2510); Algebra, label algebra-lemma-one-equation (proof line 2518); Properties, label properties-lemma-codimension-local-ring (proof line 2519). The definitions give part (1); for part (2) the local equation of the divisor and the one-equation lemma give that the local ring at the generic point has dimension one, whence the dimension function drops by one.
\end{itemize}\medskip
\ifdefined\hypertarget\hypertarget{sourcepointer-1}{}\fi
\noindent\textit{Source reference.} Original lemma 00KW: sources/stacks-chow/tags.txt, source lines 766--766. Navigation pointer for \textbackslash{}ref\{algebra-lemma-one-equation\} used inside the counted span (line 2518); the argument that consumes it is retained in the span.\par\medskip
\ifdefined\hypertarget\hypertarget{sourcepointer-2}{}\fi
\noindent\textit{Source reference.} Original definition 0BE2: sources/stacks-chow/tags.txt, source lines 14335--14335. Navigation pointer for \textbackslash{}ref\{divisors-definition-Weil-divisor\} used inside the counted span (line 2508); the argument that consumes it is retained in the span.\par\medskip
\ifdefined\hypertarget\hypertarget{sourcepointer-3}{}\fi
\noindent\textit{Source reference.} Original lemma 0AG8: sources/stacks-chow/tags.txt, source lines 13186--13186. Navigation pointer for \textbackslash{}ref\{divisors-lemma-effective-Cartier-in-points\} used inside the counted span (line 2516); the argument that consumes it is retained in the span.\par\medskip
\ifdefined\hypertarget\hypertarget{sourcepointer-4}{}\fi
\noindent\textit{Source reference.} Original lemma 02IZ: sources/stacks-chow/tags.txt, source lines 3318--3318. Navigation pointer for \textbackslash{}ref\{properties-lemma-codimension-local-ring\} used inside the counted span (line 2519); the argument that consumes it is retained in the span.\par\medskip
\ifdefined\hypertarget\hypertarget{sourcepointer-5}{}\fi
\noindent\textit{Source reference.} Original reference 02QL: sources/stacks-chow/pinned-original-source.tex, source lines 1372--1372. Navigation pointer for \textbackslash{}ref\{situation-setup\} used inside the counted span (line 2490); the argument that consumes it is retained in the span.\par\medskip
\ifdefined\hypertarget\hypertarget{sourcepointer-6}{}\fi
\noindent\textit{Source reference.} Original definition 02I9: sources/stacks-chow/tags.txt, source lines 3292--3292. Navigation pointer for \textbackslash{}ref\{topology-definition-dimension-function\} used inside the counted span (line 2510); the argument that consumes it is retained in the span.\par\medskip
\begin{lemma}
\label{lemma-divisor-delta-dimension}
Let $(S, \delta)$ be as in Situation \ref{situation-setup}.
Let $X$ be locally of finite type over $S$. Assume $X$ is
integral.
\begin{enumerate}
\item If $Z \subset X$ is an integral closed subscheme, then
the following are equivalent:
\begin{enumerate}
\item $Z$ is a prime divisor,
\item $Z$ has codimension $1$ in $X$, and
\item $\dim_\delta(Z) = \dim_\delta(X) - 1$.
\end{enumerate}
\item If $Z$ is an irreducible component of an effective Cartier
divisor on $X$, then $\dim_\delta(Z) = \dim_\delta(X) - 1$.
\end{enumerate}
\end{lemma}

\begin{proof}
Part (1) follows from the definition of a prime divisor
(Divisors, Definition \ref{divisors-definition-Weil-divisor})
and the definition of a dimension function
(Topology, Definition \ref{topology-definition-dimension-function}).
Let $\xi \in Z$ be the generic point of an irreducible component $Z$ of
an effective Cartier divisor $D \subset X$.
Then $\dim(\mathcal{O}_{D, \xi}) = 0$ and
$\mathcal{O}_{D, \xi} = \mathcal{O}_{X, \xi}/(f)$ for some
nonzerodivisor $f \in \mathcal{O}_{X, \xi}$ (Divisors,
Lemma \ref{divisors-lemma-effective-Cartier-in-points}).
Then $\dim(\mathcal{O}_{X, \xi}) = 1$ by
Algebra, Lemma \ref{algebra-lemma-one-equation}. Hence $Z$ is as in (1) by
Properties, Lemma \ref{properties-lemma-codimension-local-ring}
and the proof is complete.
\end{proof}

\end{document}
